% AY2021 材料力学 〔I〕（転記）
% 出典: 04_材料力学/original/04_材料力学/2021/2021za.pdf
% 要約: 丸棒1(直径d, 縦弾性係数E), 丸棒2(直径d, 縦弾性係数3E), 円管3(内径4d, 外径6d, 2E)
%       からなる組合せ棒。下面固定, 上面剛体円板で固定, 2本の丸棒の接合点Bに外力P。剛体板平行。
%       丸棒1は上部(A-B, 長さL/2), 丸棒2は下部(B-C, 長さL/2), 円管3は全長L。
%       (1) N1,N2,N3, (2) σ1,σ2,σ3, (3) 伸びλ1,λ2,λ3,
%       (4) Pに加え上部剛体板に中心軸まわりトルクT負荷時のねじれ角Ψ（横弾性係数は全てG）。

\section*{〔I〕}

丸棒 1（直径 $d$, 縦弾性係数 $E$）, 丸棒 2（直径 $d$, 縦弾性係数 $3E$）および
円管 3（内径 $4d$, 外径 $6d$, 縦弾性係数 $2E$）からなる組合せ棒に関する次の問い
(1)\,$\sim$\,(4) に答えよ. 組合せ棒の下面は固定されており, 上面は剛体円板で固定され,
2 本の丸棒の接合点 B に外力 $P$ が作用している. 剛体板は平行を保つものとする.

\begin{center}
% 断面図（同心円）
\begin{tikzpicture}[>=Latex,scale=0.95]
  \begin{scope}
    \clip (0,0) circle (1.8); \fill[pattern=north east lines] (-2,-2) rectangle (2,2);
    \fill[white] (0,0) circle (1.2);
  \end{scope}
  \draw (0,0) circle (1.8); \draw (0,0) circle (1.2);
  % 丸棒1/丸棒2（直径d=半径0.3, 中心）
  \fill[gray!30] (0,0) circle (0.3); \draw (0,0) circle (0.3);
  \node at (0,0.85) {\footnotesize 丸棒 1／丸棒 2};
  \draw[->] (0,0.7) -- (0,0.34);
  \node[below left] at (-1.5,-1.2) {\footnotesize 円管 3};
  \draw[->] (-1.25,-1.05) -- (-0.95,-0.85);
  \draw[->,line width=1pt] (1.4,1.6) arc (55:-35:2.1) node[midway,right]{$T$};
\end{tikzpicture}
\hspace{8mm}
% 縦断面図
\begin{tikzpicture}[>=Latex,scale=0.95]
  \draw[thick] (-2.0,0) -- (2.0,0);
  \foreach \x in {-1.9,-1.6,...,1.95} {\draw (\x,0) -- (\x-0.25,-0.25);}
  \draw[line width=2pt] (-2.0,5.0) -- (2.0,5.0);
  % 円管3（外側2壁, 全長L）
  \fill[gray!20] (-1.5,0) rectangle (-1.2,5.0);
  \fill[gray!20] (1.2,0) rectangle (1.5,5.0);
  \draw (-1.5,0) rectangle (-1.2,5.0); \draw (1.2,0) rectangle (1.5,5.0);
  % 丸棒1（上部 A-B, 直径d）
  \draw (-0.3,2.5) rectangle (0.3,5.0);
  % 丸棒2（下部 B-C, 直径d, 灰）
  \fill[gray!35] (-0.3,0) rectangle (0.3,2.5); \draw (-0.3,0) rectangle (0.3,2.5);
  \node[left] at (-0.3,4.6) {A}; \node[left] at (-0.3,2.6) {B}; \node[right] at (0.3,0.3) {C};
  \fill (0,2.5) circle (1.6pt);
  \draw[->,very thick] (0,2.5) -- (0,1.4) node[midway,right]{$P$};
  \node[right] at (0.35,3.9) {\footnotesize 丸棒 1};
  \node[right] at (0.35,1.4) {\footnotesize 丸棒 2};
  \node[right] at (1.55,4.3) {\footnotesize 円管 3};
  \draw[<->] (-1.85,0) -- (-1.85,5.0) node[midway,fill=white]{$L$};
  \draw[<->] (-0.8,0) -- (-0.8,2.5) node[midway,fill=white,font=\scriptsize]{$L/2$};
\end{tikzpicture}

\vspace{1mm}
図1
\end{center}

\begin{enumerate}[label=(\arabic*)]
  \item 丸棒 1, 2 および円管 3 に生じる内力 $N_1$, $N_2$, $N_3$ を求めよ.

  \item 丸棒 1, 2 および円管 3 に生じる応力 $\sigma_1$, $\sigma_2$, $\sigma_3$ を求めよ.

  \item 丸棒 1, 2 および円管 3 の伸び $\lambda_1$, $\lambda_2$, $\lambda_3$ を求めよ.

  \item 外力 $P$ に加え, 図に示すように上部の剛体板に中心軸周りにトルク（ねじりモーメント）
        $T$ を負荷した時のねじれ角 $\Psi$ を求めよ. なお, 丸棒 1, 丸棒 2 および円管 3 の
        横弾性係数を全て $G$ とせよ.
\end{enumerate}
